%% appendix-nsse — results of Zhao26 (plain entailment) used beyond the
%% background: satisfiability and clashes (with copy literals), cores, ordinary
%% and symbolic states of derived monitors (the monitors, derivatives and label
%% facts are in background.tex), the structure of the bounded-depth search, and
%% its constants with the closed-form estimate.

\section{Results Used from the Plain-Entailment Proof}\label{app:bg}

This appendix collects what we use from the decidability proof for plain
entailment of \citet{Zhao26} beyond \cref{sec:bg}: the satisfiability criteria
(\cref{app:bg-sat}); the cores of the
canonical hierarchy (\cref{app:bg-cores}); the ordinary and symbolic states
of derived monitors (\cref{app:bg-monitors}); the structure of the proof of
\cref{thm:bg-search} in the precise terms that later sections extend
(\cref{app:bg-search}); and the constants of its length bound
(\cref{app:bg-constants}). The results are quoted
with the hypotheses of that paper and proved there, for every arity and every
variance; we give proofs only where we extend a result. Except in
\cref{app:bg-sat}, nothing here uses a property of subtype constraints, and
every alphabet may be infinite unless it is said to be finite.

\subsection{Satisfiability}\label{app:bg-sat}

In the syntax of \citet{Zhao26} a system has no copy literals
(\cref{sec:bg-nsse}). For a covariant system $\chi$ without copy literals over a
finite set $V$, the \emph{closure} $D^{\mathrm{cov}}_\chi$ is the least reflexive and
transitive relation on $V$ such that, for every $i\in\Sigma$,
\begin{equation}\label{eq:bg-decomp}
 \frac{f(a_j)_j\le u\in\chi\qquad D^{\mathrm{cov}}_\chi(u,v)\qquad v\le f(b_j)_j\in\chi}
      {D^{\mathrm{cov}}_\chi(a_i,b_i)} .
\end{equation}
It is computed by saturation, which terminates since there are $\len V^2$
pairs. The \emph{path bounds} $\Up(\eta,u,v)$ and
$\Lo(\eta,u,v)$, for words $\eta\in\Sigma^*$ and $u,v\in V$, are defined by
$\Up(\eps,u,v)\iff\Lo(\eps,u,v)\iff D^{\mathrm{cov}}_\chi(u,v)$ and
\begin{equation}\label{eq:bg-paths}
\begin{aligned}
 \Up(i\eta,u,v)&\iff\exists z,\bar b\ \bigl[D^{\mathrm{cov}}_\chi(u,z)\land
   (z\le f(b_j)_j)\in\chi\land\Up(\eta,b_i,v)\bigr],\\
 \Lo(i\eta,u,v)&\iff\exists\bar a,z\ \bigl[(f(a_j)_j\le z)\in\chi\land
   D^{\mathrm{cov}}_\chi(z,v)\land\Lo(\eta,u,a_i)\bigr].
\end{aligned}
\end{equation}
In every solution, $\Up(\eta,u,v)$ says that the subtree of $u$ at $\eta$, when
it exists, is below $v$, and $\Lo(\eta,u,v)$ says that $u$ is below the subtree
of $v$ at $\eta$, when it exists \nssecite{Lemma ``Path rules''}. A \emph{label clash} is a
literal $f(\bar a)\le u$ or $\Top\le u$ of $\chi$, with \emph{root label} $f$ or
$\Top$, and a literal $v\le f(\bar b)$ or $v=\Bot$, with root label $f$ or
$\Bot$, such that $D^{\mathrm{cov}}_\chi(u,v)$ and the first root label is greater than the
second; that is, every such pair except two constructor literals. A
\emph{cycle clash} is a nonempty word $\eta$ and variables $a,b$ with
\begin{equation}\label{eq:bg-cycle-clash}
 \Lo(\eta,a,a),\qquad D^{\mathrm{cov}}_\chi(a,b),\qquad \Up(\eta,b,b).
\end{equation}

\begin{theorem}[Satisfiability criteria {\nssecite{Sect.\ ``Closure and Supported Path Bounds''; Theorems ``Least shape and satisfiability'', ``Covariant finite satisfiability'' and ``Regular countermodel transfer''; Lemma ``Common shape and selection''; Lemma ``Finite signed correspondence''}}]\label{thm:bg-nsse-sat}
For every arity and every variance:
\begin{enumerate}[leftmargin=*]
\item Let $\chi$ be a covariant system without copy literals. In every
  solution $\sigma$ of $\chi$, $D^{\mathrm{cov}}_\chi(u,v)$ implies $\sigma(u)\le\sigma(v)$.
  The system $\chi$ has a solution iff it has no label clash, and a finite
  solution iff it has neither a label clash nor a cycle clash. If $\chi$ has no
  label clash, then it has a cycle clash iff it has one whose word $\eta$ has
  length at most $\len V^2$.
\item For every system $\chi$ without copy literals and every variance
  $c$, $\Sat_c(\chi)\iff\Sat(\chi^\pm)$ and
  $\Sat_{\mathrm{fin},c}(\chi)\iff\Sat_{\mathrm{fin}}(\chi^\pm)$.
\item Every system without copy literals that has a solution has a regular
  solution.
\end{enumerate}
\end{theorem}

The signed system of a system without copy literals again has no copy
literals, so part (1) decides both satisfiability problems of part (2). The
extension of parts (2) and (3) to copy literals and to all three tree classes,
\cref{lem:bg-signed-sat}, serves only the complexity bounds of the part of the
paper that is not mechanized (\cref{sec:cx}); the decision procedure does not
need it.

\subsection{Cores}\label{app:bg-cores}

Let $w$ be a nonempty word with canonical hierarchy $(w_j,\ell_j)$. Every
position of $w_j$ carries a \emph{core}, which is either a cut of $w$, an
\emph{ordinary core}, or a symbolic \emph{start core}. At level $0$, position
$i$ has core $i$. At level $j+1$, the first position has a new start core
$\star_{j+1}$, and position $i\ge1$ inherits the core of the $(i-1)$st
occurrence of $\ell_j$ in $w_j$, counting occurrences from $0$; this is the
marker that closes the block $X_{i-1}$, and the block $X_i$ is read from it.
Let $\Cores_j$ be the set of ordinary cores of $w_j$.

\begin{lemma}[Cores {\nssecite{Lemma ``Cores''}}]\label{lem:bg-cores}
\leavevmode
\begin{enumerate}[leftmargin=*]
\item $\Cores_{j+1}\subseteq\Cores_j$. The start cores of $w_j$ occupy an
  initial segment of its positions, are pairwise distinct, and number at most
  $j$; the ordinary cores increase from left to right.
\item If $\Cores_j\ne\varnothing$ and $e_j=\max\Cores_j$, then
  \[
   \Cores_{j+1}=\{p : p<e_j\text{ and }\wsuf w{e_j}\pre\wsuf wp\}.
  \]
  In particular, $(p,e_j)$ is a comparison of $w$ for every
  $p\in\Cores_{j+1}$. If $\Cores_j=\varnothing$, then
  $\Cores_{j+1}=\varnothing$.
\end{enumerate}
\end{lemma}

The cut $e_j$ is the \emph{terminal core} of level $j$, the core of the last
position of $w_j$. Thus each level collects the left ends of the comparisons
whose right end is the terminal core of the level below. The proof of
\cref{thm:bg-depth} combines this lemma with \cref{lem:bg-return-gap}. Both
lemmas hold for every word, respectively every rejected word, and every reader.

\subsection{Lettered Monitors and Derivatives}\label{app:bg-monitors}

Lettered monitors (\cref{def:bg-monitor}), the lettered monitor of a reader,
which has the same nonempty holes as the reader, shifted blocks, derivatives
(\cref{def:bg-derivative}), the label facts \cref{eq:bg-label-facts} and the
desubstitution lemma (\cref{lem:bg-desub}) are stated in
\cref{sec:bg-hierarchy}. We add the structure of the states along a
hierarchy.

\paragraph{Ordinary and symbolic states}
Along a hierarchy, the states of $\mathcal M_i$ are the $N$ states of
$\mathcal M_0$, the \emph{ordinary} states, and the $i$ \emph{symbolic} states
$\star_1,\ldots,\star_i$ added by the derivatives. The ordinary states form a
closed copy: a derived transition from an ordinary state runs the previous
monitor over $\ell X$, so it stays ordinary. The \emph{entry state} of a
position $x$ of $w_i$ is the state of $\mathcal M_i$ at the cut $x$, from which
its letter is read. By the label facts and the core recursion, a position with
an ordinary core $p$ has as entry state the state of $\mathcal M_0$ at the cut
$p$ of $w$, and a position with a start core has a symbolic entry state
\nssecite{Appendix ``Details of the Bounded-Depth Compression''}.

\subsection{Structure of the Bounded-Depth Search}\label{app:bg-search}

Let $\mathcal M_0$ be a lettered monitor with $N\ge1$ states, for instance the
monitor of a reader, and let $w$ be a nonempty hole of depth $D\le d$, with
hierarchy $(w_i,\ell_i)$ and derived monitors $\mathcal M_i$. The proof of
\cref{thm:bg-search} has the five steps sketched in \cref{sec:bg-hierarchy}.

\paragraph{Step 1: derivatives}
By \cref{lem:bg-desub}, $w_i$ is a hole of $\mathcal M_i$, which has $N+i$
states.

\paragraph{Step 2: dictionary normal form}
For a word $v$ let $\alphof(v)$ be its set of letters. The level-$i$
\emph{dictionary} is $\mathcal A_i=\alphof(w_i)$, and $g_i=\len{\mathcal A_i}$.
The \emph{start count} $s_i\le i$ is the number of positions of $w_i$ with a
start core; by \cref{lem:bg-cores} they form the \emph{start prefix}
$\wpre{w_i}{s_i}$. The set $U_i$ of \emph{entry states} is the set of ordinary
states together with the symbolic states that are entry states of positions in
the start prefix, so $\len{U_i}\le N+s_i$. Every entry state occurring in $w_i$
lies in $U_i$. The \emph{protected} letters are
\[
 P_i=\{\ell_i\}\cup\alphof(\wpre{w_i}{s_i}),\qquad\len{P_i}\le s_i+1 .
\]
An unprotected letter occurs only outside the start prefix, so all its
occurrences have ordinary entry states. A comparison $(s,e)$ of $w_i$ is
\emph{left maximal} if $s\ge1$ and its predecessor letters $w_i[s-1]$ and
$w_i[e-1]$ differ. The \emph{boundary pairs} $\operatorname{BP}_i$ are the
ordered pairs $(w_i[s-1],w_i[e-1])$ over the left-maximal comparisons with
$e<\len{w_i}$, and the \emph{endpoint pairs} $\operatorname{EP}_i$ the same
pairs over those with $e=\len{w_i}$; each pair is counted once, and
$b_i=\len{\operatorname{BP}_i}$. Endpoint pairs have $\ell_i$ as second
component. A letter $a\in\mathcal A_i\setminus P_i$ is \emph{isolated} if no pair
in $\operatorname{BP}_i$ has $a$ as one component and a letter outside $P_i$ as
the other; then all unequal comparison partners of $a$, endpoint partners
included, are protected. The \emph{exchange type} of an unprotected letter $a$
consists of the following tables of $\mathcal M_i$:
\begin{itemize}[leftmargin=*]
\item the transition $\zeta\mapsto\delta_i(\zeta,a)$ and the target
  $\Tfin_i(\zeta,a)$ for every ordinary state $\zeta$;
\item for each $R\in\{\Lambda_i,\Lfin_i\}$, the \emph{self table}
  $R((\zeta,a),(\zeta',a))$ for ordinary $\zeta,\zeta'$;
\item for each such $R$, the \emph{protected rows} $R((\zeta,a),(\vartheta,z))$
  and $R((\vartheta,z),(\zeta,a))$ for ordinary $\zeta$, $\vartheta\in U_i$, and
  $z\in P_i$;
\item for each such $R$, the \emph{initial row} $R(\triangleright,(\zeta,a))$
  for ordinary $\zeta$, which is false at every level $i\ge1$.
\end{itemize}

\begin{lemma}[Exchange types {\nssecite{Lemma ``Exchange types''}}]\label{lem:bg-exchange-types}
At every level $1\le i<D\le d$ there are at most $\mathsf C(N,d)$ exchange
types, with $\mathsf C$ from \cref{eq:bg-col}.
\end{lemma}

\begin{lemma}[Comparison-preserving maps {\nssecite{Lemma ``Comparison-preserving maps''}}]\label{lem:bg-comparison-map}
Let $v$ be a nonempty word with last letter $\ell$, and $h$ a letter map that
\emph{identifies no unequal predecessor pair}: for every comparison $(s,e)$ of
$v$ with $s\ge1$, terminal ones included, $h(v[s-1])=h(v[e-1])$ implies
$v[s-1]=v[e-1]$. Then $v$ and $h(v)$ have exactly the same comparisons, and on
the letters of $v$, $h(a)=h(\ell)$ implies $a=\ell$.
\end{lemma}

A \emph{level map} at level $i$ is a map $h$ from $\mathcal A_i$ into itself. It
acts letterwise on $w_i$, and $h(w_i)$ is expanded with the original lower
markers $\ell_{i-1},\ldots,\ell_0$ to a new word $w'$. The induced maps on the
higher levels are $h_0=h$ and $h_{r+1}([X])=[h_r^*(X)]$, where $h_r^*$ applies
$h_r$ letterwise to a block.

\begin{lemma}[Shape under a level map {\nssecite{Lemma ``Shape under a level map''}}]\label{lem:bg-level-maps}
Let $h$ be a level map at level $i$ that identifies no unequal predecessor pair
of $w_i$, fixes the start prefix $\wpre{w_i}{s_i}$ letterwise, preserves the
level-$i$ state run, and makes $h(w_i)$ a hole of $\mathcal M_i$. Then $w'$ is a
hole of $\mathcal M_0$, its markers below $i$ are unchanged, and its dictionaries
below $i$ are subsets of the old ones. For every $r\ge0$,
$w'_{i+r}=h_r(w_{i+r})$, $\ell'_{i+r}=h_r(\ell_{i+r})$, and $h_r$ identifies no
unequal predecessor pair of $w_{i+r}$. Every start count is unchanged, and
exactly the same levels are unary; in particular $w'$ has depth $D$ and no
dictionary of $w'$ is larger than the corresponding one of $w$.
\end{lemma}

\begin{lemma}[Isolated exchange {\nssecite{Lemma ``Isolated exchange''}}]\label{lem:bg-isolated-exchange}
Let $1\le i<D$, and let distinct isolated letters $a,b$ of $w_i$ have the same
exchange type, where the level-zero expansion of $b$ is not longer than that
of $a$. The level map that replaces $a$ by $b$ satisfies the hypotheses of
\cref{lem:bg-level-maps}, so the resulting word is a hole of depth $D$, no
longer than $w$, with every start count unchanged and no dictionary larger;
its level-$i$ dictionary is strictly smaller.
\end{lemma}

Of two isolated letters with the same exchange type, one always has the
shorter expansion, so the lemma applies in one of the two directions.

A hole is \emph{normal} if, at every level $1\le i<D$, distinct isolated
letters have distinct exchange types, computed in its own hierarchy.

\begin{lemma}[Dictionary normal form and bounded dictionaries {\nssecite{Lemmas ``Dictionary normal form'', ``Dictionary counting and propagation'' and ``Bounded dictionaries''}}]\label{lem:bg-dictionary}
\leavevmode
\begin{enumerate}[leftmargin=*]
\item Every nonempty hole has a normal hole of the same depth, no greater
  length, the same start counts, and no larger dictionary at any level.
\item For a normal hole of depth $D\le d$ and $1\le i<D$,
  $g_i\le\mathsf A(N,d)+2b_i$. For every nonempty word and every level $i$,
  $b_i\le b_{i+1}+g_{i+1}-1$.
\item For a normal hole of depth $D\le d$, with $\mathsf A=\mathsf A(N,d)$, at
  every $1\le i<D$,
  \[
   2b_i\le(\mathsf A-1)(3^{D-1-i}-1),\qquad g_i\le1+(\mathsf A-1)3^{D-1-i}.
  \]
  Consequently every nonempty hole of depth $D\le d$ has a hole of depth $D$ and
  no greater length with $g_i\le\mathsf G(N,d)$ for every $1\le i\le D$.
\end{enumerate}
\end{lemma}

The counting in part (2) partitions $\mathcal A_i$ into the at most $d+1$
protected letters, the isolated letters, at most one per exchange type, and the
remaining letters, each of which is a component of a boundary pair with both
components unprotected. For the propagation, \cref{lem:bg-alignment} shows
that every boundary pair of level $i$ consists of the first unequal letters,
scanning from the right, of the two blocks of a boundary or endpoint pair of
level $i+1$, with $\ell_i$ standing for an exhausted block. No bound on $g_0$
is needed.

\paragraph{Step 3: the top level}

\begin{lemma}[Unary words {\nssecite{Appendix ``Details of the Bounded-Depth Compression''}}]\label{lem:bg-unary}
Let $z^m$ be a hole of a lettered monitor with state set $\mathcal Q$. Then
$z^{m'}$ is a hole for some $1\le m'\le\len{\mathcal Q}$.
\end{lemma}
\begin{proof}
Let $\zeta_k$ be the state at the cut $k$ of $z^m$, so that the label of the
cut $k\ge1$ is $(\zeta_{k-1},z)$. The states $\zeta_0,\zeta_1,\ldots$ are the
orbit of the start state under the letter $z$, so every state of the orbit
occurs first at an index below $\len{\mathcal Q}$; let $t$ be the least index
with $\zeta_t=\zeta_{m-1}$, so that $t<\len{\mathcal Q}$ and $t\le m-1$. The
prefix $z^{t+1}$ has the labels of $z^m$ at its cuts, and its last label
$(\zeta_t,z)$ is the last label of $z^m$. Every pair of cuts of a unary word
is a comparison, so a nonterminal comparison of $z^{t+1}$ is one of $z^m$ with
the same labels, and a terminal one $(s,t+1)$ has the labels of the terminal
comparison $(s,m)$ of $z^m$. Hence $z^{t+1}$ is a hole.
\end{proof}

Applied to $w_D$ in $\mathcal M_D$, the lemma gives at most $N+D$ letters.

\paragraph{Step 4: simultaneous compression}
Let $\mathcal M$ be a lettered monitor with state set $\mathcal Q$ and start
$\zeta_0$, and $\ell$ a marker. Write $\lab_0(\zeta)$ for the initial
shifted-block label: $(\zeta,\ell)$ for $\zeta\in\mathcal Q$ and $\triangleright$
for $\star$. For blocks $X_1,\ldots,X_k$ avoiding $\ell$, their \emph{summary}
is the tuple of the four functions
\[
\begin{aligned}
 \fexit_a(\zeta)&\in\mathcal Q,&&\text{the exit of the shifted block of }X_a
   \text{ read from }\zeta,\\
 \fterm_a(\zeta,\zeta'')&\in\{0,1\},&&\text{whether some label of it is related
   by }\Lfin\text{ to }(\zeta'',\ell),\\
 \feq(a,b)&\in\{0,1\},&&\text{whether }X_a=X_b,\\
 \fadm(a,b,\zeta,\zeta')&\in\{0,1\},&&\text{whether }\Lambda_\ell\text{ relates }
   (\zeta,X_a)\text{ and }(\zeta',X_b),
\end{aligned}
\]
for blocks $a,b$, entry states $\zeta,\zeta'\in\mathcal Q\sqcup\{\star\}$, and
states $\zeta''\in\mathcal Q$. It determines the restriction of $\mathcal M_\ell$
to the letters $X_1,\ldots,X_k$, together with their equalities, and it has at
most $\mathsf{col}(\len{\mathcal Q},k)$ values \cref{eq:bg-col}.

The \emph{column word} of the tuple aligns the blocks at their right ends, pads
them on the left, and reads the resulting word of $k$-tuples of optional
letters, the \emph{columns}, from left to right. It is \emph{valid} if every
track consists of absent entries followed by letters; a subsequence of a valid
column word is valid. A deterministic automaton computes the summary along the
column word. The tuple of $k$ empty blocks has
$\fexit_a(\zeta)=\delta(\zeta,\ell)$ for $\zeta\in\mathcal Q$,
$\fexit_a(\star)=\zeta_0$, $\fterm_a(\zeta,\zeta'')=
\Lfin(\lab_0(\zeta),(\zeta'',\ell))$, all blocks equal, and
$\fadm(a,b,\zeta,\zeta')=\Lambda(\lab_0(\zeta),\lab_0(\zeta'))$. Reading a column
$\gamma$ updates all components simultaneously, from the old values. If
$\gamma_a$ is absent, the components of block $a$ are unchanged. If
$\gamma_a=z$, then
\[
 \fexit'_a(\zeta)=\delta(\fexit_a(\zeta),z),\qquad
 \fterm'_a(\zeta,\zeta'')=\fterm_a(\zeta,\zeta'')\lor
   \Lfin\bigl((\fexit_a(\zeta),z),(\zeta'',\ell)\bigr).
\]
Equality becomes $\feq(a,b)\land\gamma_a=\gamma_b$. If $\gamma_a=z$ and
$\gamma_b=z'$ are both present,
\[
 \fadm'(a,b,\zeta,\zeta')=
 \Lambda\bigl((\fexit_a(\zeta),z),(\fexit_b(\zeta'),z')\bigr)
 \lor\bigl(z=z'\land\fadm(a,b,\zeta,\zeta')\bigr);
\]
if only $\gamma_a=z$ is present,
$\fadm'(a,b,\zeta,\zeta')=\Lambda((\fexit_a(\zeta),z),\lab_0(\zeta'))$; if only
$\gamma_b=z'$ is present,
$\fadm'(a,b,\zeta,\zeta')=\Lambda(\lab_0(\zeta),(\fexit_b(\zeta'),z'))$; and if
both are absent, the bit is unchanged. The update is correct because
$\Lambda_\ell$ compares positions at the same distance from the ends of two
blocks, with equal letters after them: appending a column either relates the
two new last positions, or extends the common suffix of a pair related earlier
by an equal letter.

\begin{lemma}[Column shortening {\nssecite{Appendix ``Details of the Bounded-Depth Compression''}}]\label{lem:bg-columns}
Deleting the columns between two cuts of a valid column word with equal
summaries gives a valid column word with the same final summary. Hence every
tuple of $k$ blocks avoiding $\ell$ can be replaced by subsequences, one per
block, with the same summary and all lengths below
$\mathsf{col}(\len{\mathcal Q},k)$. The replacement is injective on the
letters, preserves marker avoidance, maps holes of $\mathcal M_\ell$ over these
letters to holes, and keeps a unary level unary.
\end{lemma}

For lettered monitors $\mathcal U,\mathcal V$ with the same state set, a finite
set $E$ of letters, and a map $\rho$ from $E$ into the alphabet of $\mathcal V$,
\emph{agreement} $\mathcal U\equiv_{E,\rho}\mathcal V$ means: their start states
coincide, $\rho$ is injective on $E$, and for every $a\in E$ and every state
$\zeta$,
\begin{equation}\label{eq:bg-agreement}
 \delta_{\mathcal V}(\zeta,\rho(a))=\delta_{\mathcal U}(\zeta,a),\quad
 \Tfin_{\mathcal V}(\zeta,\rho(a))\iff\Tfin_{\mathcal U}(\zeta,a),\quad
 R_{\mathcal V}(\widehat\rho L,\widehat\rho L')\iff R_{\mathcal U}(L,L')
\end{equation}
for $R\in\{\Lambda,\Lfin\}$ and all labels $L,L'$ whose letters lie in $E$,
the start label included; here $\widehat\rho(\triangleright)=\triangleright$ and
$\widehat\rho(\zeta,a)=(\zeta,\rho(a))$, and states are not renamed. Agreement
preserves runs, labels up to $\widehat\rho$, and holes of words over $E$, and
agreements compose.

\begin{lemma}[Agreement under derivation {\nssecite{Lemma ``Agreement under derivation''}}]\label{lem:bg-agreement}
If $\mathcal U\equiv_{E,\rho}\mathcal V$ and $\ell\in E$, then for every finite
family $F$ of blocks over $E$ avoiding $\ell$, the letterwise map $\rho^*$ gives
$\mathcal U_\ell\equiv_{F,\rho^*}\mathcal V_{\rho(\ell)}$, and the image blocks
avoid $\rho(\ell)$.
\end{lemma}

For markers $\boldsymbol\ell=(\ell_0,\ldots,\ell_{j-1})$ write
$\mathcal M^{\boldsymbol\ell}_0=\mathcal M_0$ and
$\mathcal M^{\boldsymbol\ell}_{h+1}=(\mathcal M^{\boldsymbol\ell}_h)_{\ell_h}$. A
finite family of level-$j$ letters is \emph{valid} for the markers if every
block avoids the marker of its level and its letters are, recursively, valid
letters of the level below; $\Exp_{\boldsymbol\ell,j}(a)$ is the level-zero
word obtained by expanding a valid level-$j$ letter $a$, each block followed by
its marker.

\begin{lemma}[Supported simultaneous compression {\nssecite{Lemma ``Supported simultaneous compression''}}]\label{lem:bg-compression}
Fix a nonempty word $w$ with hierarchy markers $\boldsymbol\ell$ and
dictionaries $\mathcal A_h$, and $j\le d$. Suppose $\len{\mathcal A_h}\le g$ for
$1\le h\le j$, and put $\kappa=\mathsf{col}(N+d+1,g)$. For every set
$E_j\subseteq\mathcal A_j$ there are new markers $\widetilde{\boldsymbol\ell}$
and a map $\rho_j$, injective on $E_j$, such that
\[
 \mathcal M^{\boldsymbol\ell}_j\equiv_{E_j,\rho_j}
 \mathcal M^{\widetilde{\boldsymbol\ell}}_j,
 \qquad
 \len{\Exp_{\widetilde{\boldsymbol\ell},j}(\rho_j(a))}\le\kappa^j\quad(a\in E_j),
\]
and the image family is valid for the new markers. The recursion keeps at each
level $h<j$ the support
\[
 E_h=\{\ell_h\}\cup\bigcup_{X\in E_{h+1}}\alphof(\tau_h(X))\ \subseteq\ \mathcal A_h,
\]
where $\tau_h(X)$ is the subsequence of $X$ produced by \cref{lem:bg-columns};
at level zero no cardinality bound is needed. Agreement holds at every state,
not only at the states visited by $w$.
\end{lemma}

\paragraph{Step 5: the bound}
Normalize the given hole by \cref{lem:bg-dictionary}, so that $g_i\le
\mathsf G=\mathsf G(N,d)$ for $1\le i\le D$, and keep it fixed as the original
hierarchy; put $\kappa=\mathsf{col}(N+d+1,\mathsf G)$. \Cref{lem:bg-unary} gives
a hole $z^m$ of $\mathcal M_D$ with $1\le m\le N+D$. Compressing $E_D=\{z\}$ by
\cref{lem:bg-compression}, agreement makes $\rho_D(z)^m$ a hole of the new
level-$D$ monitor, and validity allows its exact expansion by
\cref{lem:bg-desub} to a nonempty hole of $\mathcal M_0$. Cutting that expansion
recovers the unary word $\rho_D(z)^m$ at level $D$, so the depth is at most $D$,
and the length is at most
\[
 m\kappa^D\le(N+D)\kappa^D\le(N+d+1)\kappa^d=\holeBound(N,d).
\]
No step uses finiteness of the alphabet, and the alphabet size does not enter
the bound.

\subsection{Constants}\label{app:bg-constants}

For integers $N\ge1$, $d\ge0$, and $k\ge0$ put
\begin{equation}\label{eq:bg-col}
\begin{aligned}
 \mathsf{col}(N,k)&=N^{(N+1)k}\cdot2^{N(N+1)k}\cdot2^{k^2}\cdot2^{(N+1)^2k^2},
 &\mathsf C(N,d)&=N^N\,2^{2N^2+4N(N+d)(d+1)+3N},\\
 \mathsf A(N,d)&=d+1+\mathsf C(N,d),
 &\mathsf G(N,d)&=2+\mathsf A(N,d)\,3^d,
\end{aligned}
\end{equation}
and
\begin{equation}\label{eq:bg-len}
 \holeBound(N,d)=(N+d+1)\cdot\mathsf{col}\bigl(N+d+1,\mathsf G(N,d)\bigr)^d .
\end{equation}
All five functions are nondecreasing in each argument. Their meaning is as
follows; $D\le d$ is the depth of the hole being shortened.
\begin{itemize}[leftmargin=*]
\item $\mathsf{col}(N,k)$ bounds the number of summary values of a $k$-tuple of
  blocks for a lettered monitor with $N$ states. It is the product of the
  numbers of values of the summary fields:
  \begin{center}
  \begin{tabular}{@{}ll@{}}
  \toprule
  Factor & Summary field\\
  \midrule
  $N^{(N+1)k}$ & exits: $k$ blocks, $N+1$ entry states\\
  $2^{N(N+1)k}$ & terminal-comparison bits, also for $N$ final-marker states\\
  $2^{k^2}$ & equality bits for ordered pairs of blocks\\
  $2^{(N+1)^2k^2}$ & derived-admission bits for pairs of blocks and entry states\\
  \bottomrule
  \end{tabular}
  \end{center}
\item $\mathsf C(N,d)$ bounds the number of exchange types at a level $1\le i<D$
  (\cref{lem:bg-exchange-types}). The transition table of the ordinary states
  has $N^N$ values; the target table uses $N$ bits, the two self tables $2N^2$
  bits, the protected rows $4N\len{U_i}\len{P_i}$ bits, where $\len{P_i}\le d+1$
  and $\len{U_i}\le N+d$, and the initial rows $2N$ bits.
\item $\mathsf A(N,d)$ bounds the protected letters, at most $d+1$, together with
  one isolated letter for each exchange type.
\item $\mathsf G(N,d)$ bounds the dictionary at every level $1\le i\le D$ of a
  suitable normal hole (\cref{lem:bg-dictionary}).
\item $\holeBound(N,d)$ is the length bound. The top level of the shortened hole
  has at most $N+D\le N+d+1$ letters, and each of the at most $d$ expansion
  steps multiplies the length by at most $\mathsf{col}(N+d+1,\mathsf G(N,d))$,
  which bounds the summary values of the level-$j$ monitors, with $N+j$ states,
  for every $j\le d$. For $d=0$ the bound is $N+1$: a hole of depth $0$ is unary
  and is its own top level.
\end{itemize}

\begin{lemma}[Elementary estimate {\nssecite{Lemma ``Elementary estimate''}}]\label{lem:bg-estimate}
Put $\Xi(N,d)=3N^2+4N(N+d)(d+1)+3N+3d+4$. Then
$\mathsf G(N,d)\le2^{\Xi(N,d)}$ and $\holeBound(N,d)\le2^{2^{5\Xi(N,d)+3}}$. For
$N\ge1$ and $d_N=2N(4N-1)-1$,
\[
 \holeBound(N,d_N)\le2^{2^{1280(N+1)^5}} .
\]
\end{lemma}
\begin{proof}
Write $\Xi=\Xi(N,d)$ and $\Xi_0=3N^2+4N(N+d)(d+1)+3N$. From $N\le2^N$ we get
$\mathsf C(N,d)\le2^{\Xi_0}$; a sum of two quantities bounded by $2^p$ and $2^q$
is at most $2^{p+q+1}$, and $d+1\le2^{d+1}$, so $\mathsf A(N,d)\le2^{\Xi_0+d+2}$;
with $3^d\le2^{2d}$ this gives $\mathsf G(N,d)\le2^{\Xi_0+3d+3}\le2^\Xi$. Put
$M=N+d+1$, $g=\mathsf G(N,d)$, and $P=2^\Xi$. Then $M+1\le\Xi$, so
$M,M+1,d,g\le P$. Replacing $M$ by $2^M$ in the first factor of
\cref{eq:bg-col} yields
$\mathsf{col}(M,g)\le2^{2M(M+1)g+g^2+(M+1)^2g^2}\le2^{4P^4}=2^{2^{4\Xi+2}}$.
Therefore $\holeBound(N,d)\le2^{\Xi+d2^{4\Xi+2}}\le2^{2^{5\Xi+3}}$, because
$d\le2^\Xi$ and $\Xi\le2^{5\Xi+2}$. For $d_N=8N^2-2N-1$ we have $d_N$,
$N+d_N$, $d_N+1\le8N^2$, so $4N(N+d_N)(d_N+1)\le256N^5$ and
$\Xi(N,d_N)\le256N^5+27N^2+3N+4$. Hence $5\Xi(N,d_N)+3\le1280N^5+135N^2+15N+23
\le1280(N+1)^5$.
\end{proof}
